% SecondWave — deck de présentation (10 slides)
% Compilation :  tectonic DECK.tex
\documentclass[aspectratio=169,11pt]{beamer}

\usetheme{default}
\usecolortheme{default}
\setbeamertemplate{navigation symbols}{}
\usepackage[T1]{fontenc}
\usepackage{lmodern}
\usepackage[french]{babel}
\usepackage{booktabs}
\usepackage{tabularx}
\usepackage{xcolor}

\definecolor{accent}{HTML}{0D6E6E}
\definecolor{ink}{HTML}{16202A}
\definecolor{ink2}{HTML}{56646F}
\definecolor{bad}{HTML}{A33131}
\definecolor{good}{HTML}{2C6E4C}
\definecolor{shade}{HTML}{EFF3F4}

\setbeamercolor{normal text}{fg=ink}
\setbeamercolor{frametitle}{fg=accent}
\setbeamercolor{itemize item}{fg=accent}
\setbeamercolor{itemize subitem}{fg=ink2}
\setbeamerfont{frametitle}{size=\large,series=\bfseries}
\setbeamertemplate{frametitle}{\vspace{6pt}\insertframetitle\par\vspace{-4pt}
  {\color{accent}\rule{\textwidth}{0.6pt}}}
\setbeamertemplate{itemize item}{\textbf{--}}
\setbeamertemplate{itemize subitem}{\textperiodcentered}
\setbeamertemplate{footline}{\hfill{\tiny\color{ink2}
  SecondWave \textperiodcentered\ BFT-PARIS-26 \textperiodcentered\ \insertframenumber/10}\hspace{10pt}\vspace{4pt}}

\newcommand{\tc}[1]{{\small\texttt{#1}}}

\begin{document}

% —————————————————————————————————— 1
\begin{frame}[plain]
\vspace{1.1cm}
\begin{center}
{\Huge\bfseries SecondWave}\\[10pt]
{\large Un marché secondaire pour les positions de vault verrouillées,}\\
{\large avec un analyste de risque intégré}\\[26pt]
{\color{ink2}\small
Track 2 --- Lending Protocol \textperiodcentered\ flavour Loaded\\
Devnet public \textperiodcentered\ rippled 3.4.0-rc5 \textperiodcentered\ \tc{xrpl.js@5.2.0-beta.1}\\[8pt]
\textbf{équipe BFT-PARIS-26}}
\end{center}
\end{frame}

% —————————————————————————————————— 2
\begin{frame}{Le capital est enfermé, et le risque bouge quand même}
\vspace{4pt}
Un vault \textbf{closed-ended} verrouille les dépôts pendant toute la phase Investment.
\tc{VaultWithdraw} $\rightarrow$ \tc{tecTOO\_SOON}. Un déposant qui a besoin de son argent
n'a \textbf{aucune issue} jusqu'à la \tc{RedemptionDate}.
\vspace{10pt}

\begin{center}
\begin{tabular}{@{}lll@{}}
\toprule
\textbf{Subscription} & \textbf{Investment} & \textbf{Redemption}\\
\midrule
dépôt et retrait libres & \color{bad}\textbf{tout est bloqué} & retrait du capital + rendement\\
prêter est interdit & on prête & plus de nouveau prêt\\
\bottomrule
\end{tabular}
\end{center}
\vspace{10pt}

Et pendant ce verrouillage, \textbf{le risque se dégrade sans que le prix bouge}.
\tc{AssetsTotal} ne change qu'au moment où le courtier déclare un impairment --- et
\textbf{rien ne l'y oblige}.
\vspace{4pt}

{\color{bad}\textbf{Un prêt défaillable depuis des heures laisse la NAV parfaitement intacte.}}
\end{frame}

% —————————————————————————————————— 3
\begin{frame}{Les parts sont un jeton : la sortie existe, elle n'est pas outillée}
\vspace{6pt}
Les parts d'un vault sont un \textbf{MPT} (XLS-33). Elles se transfèrent. On outille
la cession en trois briques :
\vspace{8pt}
\begin{itemize}
\setlength\itemsep{7pt}
\item \textbf{Un carnet d'offres} --- n'affiche à un acheteur que ce qu'il peut
      \emph{réellement} recevoir, et dit \emph{pourquoi} quand une offre lui est fermée.
\item \textbf{Un règlement atomique} --- parts contre prix, ensemble ou pas du tout,
      avec preflight, reconstruction du résultat et réconciliation des soldes.
\item \textbf{Un analyste} --- le vendeur brade toujours. Toute la question est
      \emph{pourquoi}.
\end{itemize}
\vspace{10pt}

\begin{center}
\begin{tabular}{@{}ll@{}}
\toprule
\color{good}\textbf{Décote de liquidité} & \color{bad}\textbf{Décote de détresse}\\
\midrule
le fonds va bien, le vendeur a besoin de cash & le fonds va mal, ça ne se voit pas encore\\
\color{good}une opportunité & \color{bad}un piège\\
\multicolumn{2}{c}{\textbf{prix affiché : identique}}\\
\bottomrule
\end{tabular}
\end{center}
\end{frame}

% —————————————————————————————————— 4
\begin{frame}{Le flux on-chain --- six normes XLS combinées}
\vspace{2pt}
{\small
\setlength\leftmargini{2.4cm}
\begin{itemize}
\setlength\itemsep{5pt}
\item[\textbf{XLS-70+80}] \tc{CredentialCreate} $\rightarrow$ \tc{CredentialAccept} $\rightarrow$ \tc{PermissionedDomainSet}
\item[\textbf{XLS-65}] \tc{VaultCreate} \emph{(privé, closed-ended)} $\rightarrow$ \tc{VaultDeposit}
\item[\textbf{XLS-66}] \tc{LoanBrokerSet} + \tc{CoverDeposit} $\rightarrow$ \tc{LoanSet} $\rightarrow$ \tc{LoanPay} / \tc{LoanManage}
\item[] {\color{bad}\textbf{$\Downarrow$ bascule en Investment --- les retraits sont bloqués}}
\item[\textbf{XLS-33+56}] le \tc{Batch} \tc{tfAllOrNothing} en trois jambes :\\[2pt]
      \tc{MPTokenAuthorize} $\rightarrow$ \tc{Payment} des parts $\rightarrow$ \tc{Payment} du prix
\end{itemize}}
\vspace{8pt}

\textbf{18 types de transaction} employés, chacun avec un lien on-chain vérifié dans le README.
\vspace{6pt}

{\small\color{ink2}Le gate du domaine s'applique partout : un non-membre reçoit \tc{tecNO\_AUTH} au
dépôt, au transfert direct, et même au dénouement d'un escrow.}
\end{frame}

% —————————————————————————————————— 5
\begin{frame}[plain]
\vspace{2.2cm}
\begin{center}
{\Huge\bfseries\color{accent}Démonstration}\\[16pt]
{\large De l'offre au règlement vérifié}\\[20pt]
{\color{ink2}\tc{npm run cli vaults} \quad\textperiodcentered\quad \tc{npm run cli book}
\quad\textperiodcentered\quad \tc{npm run cli buy}}
\end{center}
\end{frame}

% —————————————————————————————————— 6
\begin{frame}{Friction 1 --- un \texttt{Batch} ne dit pas s'il a produit un effet}
\vspace{4pt}
{\small\color{ink2}Réponse complète du serveur sur un \tc{Batch} qui n'a \emph{rien} fait :}
\vspace{3pt}

\colorbox{shade}{\parbox{\textwidth}{\vspace{3pt}\small\texttt{
accepted: true~~|~~applied: true~~|~~engine\_result: tesSUCCESS\\
engine\_result\_message: "The transaction was applied."\\
meta.AffectedNodes: [ 1 noeud ]~~$\leftarrow$ la ponction de frais}\vspace{3pt}}}
\vspace{8pt}

Sur \textbf{huit} \tc{Batch} rejoués --- cinq ayant déplacé des parts, trois n'ayant rien
fait --- \textbf{une seule combinaison apparaît}. Succès et non-effet sont indiscernables.
Et \tc{simulate} répond \tc{notImpl} sur un \tc{Batch} : \textbf{ni pré-jouer, ni constater}.
\vspace{8pt}

{\color{accent}\textbf{Ce qu'on demande}} --- le \tc{BatchExecutions} que XLS-56 spécifie déjà.
À défaut, \textbf{l'index de la jambe en échec et son code}. Un entier et un code.
\vspace{4pt}

{\small\color{ink2}L'amendement est encore en \tc{BatchV1\_1} sur une \emph{release candidate} :
c'est maintenant que ça coûte le moins cher.}
\end{frame}

% —————————————————————————————————— 7
\begin{frame}{Friction 2 --- trois drapeaux sur quatre livrent sans encaisser}
\vspace{6pt}
{\small\color{ink2}Le même échange, jambe du prix volontairement impayable, rejoué sous chaque drapeau :}
\vspace{6pt}

\begin{center}
\begin{tabularx}{\textwidth}{@{}Xlll@{}}
\toprule
\textbf{Drapeau} & \textbf{Annoncé} & \textbf{Parts déplacées} & \textbf{Prix payé}\\
\midrule
\tc{tfAllOrNothing} & \tc{tesSUCCESS} & \textbf{0} & 0\\
\tc{tfOnlyOne} \textperiodcentered\ \tc{tfUntilFailure} \textperiodcentered\ \tc{tfIndependent}
  & \tc{tesSUCCESS} & {\color{bad}\textbf{300 000}} & {\color{bad}\textbf{0}}\\
\bottomrule
\end{tabularx}
\end{center}
\vspace{10pt}

L'oubli est impossible --- \tc{Flags: 0} est refusé. Le risque, c'est \textbf{le mauvais
choix}, fait en connaissance de cause. Et combiné à la friction 1, il est
\textbf{indétectable} : même code, même métadonnée, et l'actif est parti.
\vspace{8pt}

{\color{accent}\textbf{Ce qu'on demande}} --- une mise en garde dans XLS-56 pour les échanges
de valeur, et un helper de swap dans les SDK qui \emph{impose} \tc{tfAllOrNothing}.
\end{frame}

% —————————————————————————————————— 8
\begin{frame}{Friction 3 --- aucune primitive de lecture côté XLS-66}
\vspace{6pt}
\tc{vault\_info} existe. \textbf{Ni \tc{loan\_info}, ni \tc{loan\_broker\_info}, ni la liste
des détenteurs de parts.}
\vspace{6pt}

\colorbox{shade}{\parbox{\textwidth}{\vspace{3pt}\small\texttt{
vault\_info(vaultId)\\
\phantom{xx}$\rightarrow$ account\_objects(pseudo-compte du \textbf{vault}, type: 'loan\_broker')\\
\phantom{xxxx}$\rightarrow$ account\_objects(pseudo-compte du \textbf{courtier}, type: 'loan')}\vspace{3pt}}}
\vspace{8pt}

Et pour savoir qui détient les parts : \textbf{rejouer tout l'historique} du pseudo-compte
et reconstruire les soldes depuis la métadonnée. C'est la partie du code qui nous a demandé
le plus de travail --- et elle ne reconstitue que ce que le serveur sait déjà.
\vspace{8pt}

{\color{accent}\textbf{Ce qu'on demande}} --- \tc{loan\_info} et \tc{loan\_broker\_info} sur le
modèle de \tc{vault\_info}, et une commande qui liste les détenteurs d'une émission MPT.
\end{frame}

% —————————————————————————————————— 9
\begin{frame}{Ce qu'on a mesuré}
\vspace{4pt}
\begin{center}
{\small\begin{tabularx}{\textwidth}{@{}Xl@{}}
\toprule
Cas exécutés sur le Devnet public & \textbf{$\sim$420} sur une trentaine de vaults\\
Coût d'un échange complet & \textbf{200 drops} --- 0,0152 \% du prix\\
Décote de marché observée & \textbf{12 à 14 \%} selon le vault\\
\multicolumn{2}{@{}l}{\quad$\rightarrow$ \textbf{le rail coûte 800 fois moins cher que la décote}}\\
\midrule
Charger l'état complet d'un vault & 4 appels RPC, $\sim$700 ms\\
Bug bloquant trouvé & \tc{LoanSet} inutilisable --- \textbf{corrigé en \tc{5.2.0-beta.1}}\\
\bottomrule
\end{tabularx}}
\end{center}
\vspace{6pt}

{\color{accent}\textbf{Et ce qui nous a rassurés}}
\vspace{2pt}
{\small\begin{itemize}
\setlength\itemsep{2pt}
\item \textbf{Personne ne peut enfermer un déposant} : les dates ne sont pas modifiables après création.
\item \textbf{Le pseudo-compte est étanche} --- paiement, escrow, chèque : \tc{tecNO\_PERMISSION}.
\item \textbf{Aucune fuite d'arrondi} sur 145 aller-retours en XRP et en IOU.
\item \tc{tfAllOrNothing} \textbf{tient} dans tous les cas d'échec testés.
\end{itemize}}
\end{frame}

% —————————————————————————————————— 10
\begin{frame}{Ce qui reste ouvert}
\vspace{8pt}
{\color{accent}\Large\textbf{Le carnet d'ordres natif connaîtra-t-il les Permissioned Domains ?}}
\vspace{10pt}

\tc{OfferCreate} avec un MPT répond \tc{temDISABLED} --- pas \tc{temMALFORMED}. Le code
existe, l'amendement dort. C'est pour ça que notre carnet est hors chaîne.
\vspace{8pt}

Mais un carnet d'ordres est \textbf{sans permission} : il apparie deux ordres sans savoir qui
les a posés. Nos parts sont gatées. Un DEX qui accepterait les MPT sans connaître les domaines
apparierait des ordres dont le règlement serait ensuite refusé.
\vspace{10pt}

\begin{center}
\begin{tabular}{@{}ll@{}}
\toprule
\textbf{Si le DEX connaît les domaines} & notre carnet devient inutile\\
\textbf{S'il ne les connaît pas} & un marché P2P reste la seule voie pour un vault privé\\
\bottomrule
\end{tabular}
\end{center}
\vspace{8pt}

{\small\color{ink2}Dans les deux cas, l'analyste reste : un carnet affiche un prix, il ne dit
pas si c'est une aubaine ou un piège.}
\end{frame}

\end{document}
